\section{Discussion}
\label{sec:discussion}

\subsection{Main findings}
The central result is that numerical reference discrepancies can be smaller than, comparable to, or larger than selected surrogate errors depending on the reference pair and metric. Our reference-aware temporal-alignment procedure establishes a shared output timeline without modifying the underlying learning task, preventing solver-specific natural time steps from corrupting performance evaluation.

For the selected seed-$18$ Hydrostatic direct models, ConvLSTM has the lowest global relative $L_2$ ($0.203$), followed by U-FNO ($0.224$) and F-FNO ($0.230$); the paired scenario bootstrap gives the same ordering for the mean per-scenario metric. This is a fixed-checkpoint comparison rather than a population-level architecture ranking. The windowed models solve an easier first-frame-conditioned task and are not like-for-like direct forecasts.

\subsection{Reference choice and numerical comparison}
At identical requested times, the global RMSE gaps are $0.00171$ for Hydrostatic--MUSCL-HR, $0.00305$ for MUSCL-HR--Boussinesq, and $0.00374$ for Hydrostatic--Boussinesq. The six seed-$18$ cross-reference ratios range from $1.10$ to $1.72$ and all exceed one, so no tested emulator bridges its raw solver gap in this metric. These ratios remain benchmark-relative effects rather than evidence of physical convergence.

The current R6 GeoClaw comparison uses three fixed representative test cases, the same $192 \times 192$ buffered computation, the same $64 \times 64$ published fields, and the same 50 requested times as the surrogate benchmark. Mean trajectory relative $L_2$ is $0.2599$ for Hydrostatic and $0.1256$ for MUSCL-HR on this small diagnostic set. These values are descriptive inter-code comparisons, not physical-error estimates or a dataset-wide ranking.

The new comparator confirms exact frozen inputs, the initial state, the publication mapping, and the requested output times before field differences are calculated \citep{leveque2011tsunami,berger2011geoclaw}. GeoClaw uses single-level open extrapolation, whereas the in-house shallow-water references use radiation-plus-sponge boundaries; the implementations also differ in fluxes, reconstruction and limiting, topographic-source treatment, time integration, and multidimensional correction terms. The remaining difference is therefore interpreted as a combined inter-code discrepancy, not as evidence that either solver is closer to the continuum solution or to physical observations. The earlier $96$-grid ablation remains historical evidence and is not used to support the current R6 results.

\subsection{Model and transfer behavior}
ConvLSTM has the lowest selected Hydrostatic direct error, while F-FNO uses far fewer parameters than FNO and U-FNO is also competitive. Performance does not strictly scale with spectral capacity: increasing retained modes from 14 to 20 gives no improvement over the base FNO, and the eight-mode variant is slightly better. ConvLSTM has the lowest low- and mid-frequency spectral errors among the three diagnostic models, whereas F-FNO improves the low-frequency error relative to FNO but has larger mid- and high-frequency errors.

On ten fully wet GEBCO-derived crops independently rescaled to the training depth range and paired with synthetic sources, direct F-FNO has lower error than direct FNO. This supports morphology transfer under the benchmark construction, not real-event, physical-depth, shoreline, or wet--dry generalization. Rough-source, resolution-transfer, and wet--dry stress tests expose additional limits.

\subsection{Uncertainty and practical scope}
The seven-member FNO ensemble is seed-stable, with member mean relative $L_2$ ranging from $0.257$ to $0.273$ and standard deviation $0.0055$, but it is under-dispersed. Validation-fitted scalar inflation improves test $90\%$ marginal coverage from $0.682$ to $0.910$, yet the high-strength subgroup reaches only $0.796$ and the error--uncertainty correlation is $0.496$. Because coverage metrics pool spatially and temporally correlated pixels, they serve as interval-scale reliability indicators rather than event-level probability calibrators.

In summary, the study supports a controlled finite-horizon surrogate comparison and exposes clear architecture-, reference-, and subgroup-dependent error behavior. It does not make an operational speed claim; deployment would additionally require validated dimensional scaling, surveyed bathymetry, wet--dry coastal physics, and calibrated hazard-scale uncertainty quantification.
